\chapter{Теоретические основы и проблема циклических дедлоков}
\label{chap:theory}

\section{Характеристики дорожного графа и мезоскопический подход}
\label{sec:mesoscopic_approach}

Мезоскопическое моделирование транспортных потоков занимает промежуточное положение между макроскопической гидродинамикой (модели LWR) и дискретной микросимуляцией individual-level кинематики. В разработанном движке каждый агент представляет собой отдельный объект с индивидуальным маршрутом, однако движение по элементам дорожной сети регулируется пространственными очередями ограниченной емкости и агрегированными функциями задержек (BPR).

Для экспериментальных исследований выделена область дорожного графа города Казани в географических координатах:
\begin{equation}
\text{BBox} = [48.46^\circ\,\text{E},\; 55.63^\circ\,\text{N},\; 49.43^\circ\,\text{E},\; 55.93^\circ\,\text{N}].
\label{eq:kazan_bbox}
\end{equation}

После обработки картографических данныхOpenStreetMap в модуле \texttt{data-processor} сформирован ориентированный реберно-ориентированный (\emph{edge-based}) граф дорожной сети. Итоговые характеристики графа приведены в таблице~\ref{tab:graph_stats}.

\begin{table}[h!]
\centering
\caption{Структурные параметры ориентированного графа сети г.~Казани}
\label{tab:graph_stats}
\begin{tabular}{|l|c|}
\hline
\textbf{Параметр графа} & \textbf{Значение} \\ \hline
Координаты области (BBox) & $[48.46, 55.63, 49.43, 55.93]$ \\ \hline
Общее число физических ребер ($N_{\text{edges}}$) & $160\,689$ \\ \hline
Доля коротких ребер (длина $L < 50$ м) & $43.0\%$ ($\approx 69\,096$ ребер) \\ \hline
Доля микроребер (длина $L < 25$ м) & $20.0\%$ ($\approx 32\,137$ ребер) \\ \hline
Порог плотности квалификации затора ($J_{\text{thresh}}$) & $70.0\%$ \\ \hline
\end{tabular}
\end{table}

\section{Анализ причин образования циклических блокировок (grid-locks)}
\label{sec:deadlock_analysis}

Особенностью высокодетализированного графа является сохранение геометрии сложных перекрестков и круговых развязок, что приводит к появлению фрагментов с высокой концентрацией ребер малой протяженности ($L < 25$ м). При вместимости одного полосового участка в 2--4 автомобиля максимальная емкость ребра $C_e$ достигается за 2--5 секунд физического времени.

Пусть $C = (e_1, e_2, \dots, e_k, e_1)$ — замкнутый контур (цикл) в ориентированном графе. Дедлок возникает при выполнении трех условий Коффмана:
\begin{enumerate}
    \item \textbf{Взаимное исключение}: Каждая физическая ячейка на ребре $e_i$ может быть занята только одним агентом.
    \item \textbf{Удержание и ожидание}: Агент в голове очереди ребра $e_i$ ожидает освобождения места на следующем ребре $e_{i+1}$, удерживая позицию на $e_i$.
    \item \textbf{Отсутствие перераспределения}: Агент не может быть принудительно удален из очереди без внешнего управляющего механизма.
\end{enumerate}

При достижении 100\% заполнения всех ребер входящих в контур $C$, мгновенные скорости всех агентов обращаются в ноль ($v = 0$ м/с), а затор распространяется каскадно на примыкающие улицы.

\section{Сравнительный анализ методов компенсации дедлоков}
\label{sec:unblocking_methods}

В современных симуляторах применяются три основных подхода к устранению замкнутых заторов:

\begin{enumerate}
    \item \textbf{Мгновенная телепортация/деспавн (Cities: Skylines)}: Агент, находящийся в состоянии простоя свыше порогового времени, обнуляется и перемещается напрямую в конечную точку маршрута. Данный метод устраняет затор, однако приводит к искажению физических метрик времени поездки (TTI) и искусственному уничтожению трафика.
    
    \item \textbf{Виртуальный буфер SUMO MESO (Дискретный таймер)}: При простое свыше $T_{\text{teleport}} = 300$ с агент удаляется с физического ребра во внешний буфер заторов. В буфере агент с заданной минимальной скоростью $v_{\text{virt}} = 2.0$ м/с виртуально проезжает ребро и совершает попытки вернуться на следующее свободное ребро маршрута. Недостаток подхода заключается в том, что при переходе агента на короткое соседнее ребро (15--20 м) таймер ожидания сбрасывается в 0, вследствие чего на графах с 43\% коротких ребер время простоя агент вырастает до неограниченных величин.
    
    \item \textbf{Разработанный метод с непрерывным инвариантом таймера (\texttt{conttimer})}: Таймер простоя $t_{\text{wait}}$ не сбрасывается при смене ребра, если скорость движения остается ниже порога $v_{\text{jam}} < 0.1$ м/с. Накапливаемый инвариант гарантирует перемещение агента в виртуальный буфер при достижении предельного интервала $T_{\text{teleport}}$ (130 с или 300 с) независимо от числа пройденных коротких ребер.
\end{enumerate}
